\subsection{Quantitative Results}

This section presents the main quantitative results of the model and connects them to the theoretical predictions.

\subsubsection{Result 1: Stablecoin Demand and Misalignment}

Figure 1 shows stablecoin demand $q_s$ as a function of exchange rate misalignment $\theta$. Stablecoin usage increases monotonically with $\theta$.

This reflects the increasing demand for foreign currency as misalignment grows. Since stablecoins offer lower transaction costs, they capture a growing share of foreign currency demand.

\subsubsection{Result 2: Crisis Probability}

Figure 2 plots the probability of crisis as a function of $\theta$ under two scenarios: (i) cash-only economy and (ii) economy with stablecoins.

The presence of stablecoins increases crisis probability for all levels of misalignment. The difference is modest at low $\theta$ but becomes more pronounced as misalignment increases.

This result reflects the information channel: stablecoins improve signal precision, facilitating coordination and increasing the likelihood of speculative attacks.

\subsubsection{Result 3: Welfare}

Figure 3 shows welfare as a function of $\theta$. Welfare exhibits a non-monotonic relationship with misalignment.

At low levels of $\theta$, welfare is higher in the presence of stablecoins due to lower transaction costs and improved access to foreign currency. At higher levels of $\theta$, welfare is lower due to increased crisis probability.

\subsubsection{Result 4: Welfare Difference}

Figure 4 plots the difference in welfare between the stablecoin and cash-only economies.

The welfare difference is positive at low misalignment and negative at high misalignment, crossing zero at an intermediate threshold. This confirms the theoretical prediction of a welfare tradeoff.

\subsubsection{Magnitude and Realism}

Importantly, the magnitudes of the effects are moderate. Crisis probabilities remain well below one, and welfare changes are economically meaningful but not extreme.

This ensures that the results are consistent with realistic macro-financial environments rather than driven by extreme parameter choices.

\subsubsection{Link to Theory}

The quantitative results closely match the theoretical predictions:

\begin{itemize}
\item Increasing $q_s$ raises signal precision (Proposition 2)
\item Higher precision increases coordination (Propositions 5--7)
\item Welfare reflects the tradeoff between efficiency and fragility (Proposition 8)
\end{itemize}

Thus, the simulations provide a quantitative validation of the theoretical mechanism.
